\documentclass[11pt]{article}
% ======================================================================
%  Preambule commun - Sujets Bac Serie A (Mathematiques)
%  Style sobre : noir et blanc, sans couleurs, sans filigrane,
%  sans en-tetes ni pieds de page decoratifs.
% ======================================================================
\usepackage[utf8]{inputenc}
\usepackage[T1]{fontenc}
\usepackage{lmodern}
\usepackage{textcomp}
\usepackage{amsmath,amssymb}
\usepackage{enumitem}
\usepackage{array}
\usepackage[a4paper,margin=2.3cm]{geometry}
\usepackage{fancyhdr}
\usepackage{tikz}

% Symbole degre robuste + justification souple
\DeclareUnicodeCharacter{00B0}{\ensuremath{^\circ}}
\tolerance=2000
\emergencystretch=2em
\frenchspacing

% En-tete / pied de page minimalistes (numero de page seul)
\pagestyle{fancy}
\fancyhf{}
\renewcommand{\headrulewidth}{0pt}
\fancyfoot[C]{\thepage}

% Macros mathematiques
\newcommand{\R}{\mathbb{R}}
\newcommand{\N}{\mathbb{N}}
\newcommand{\Z}{\mathbb{Z}}
\newcommand{\vi}{\vec{\imath}}
\newcommand{\vj}{\vec{\jmath}}

% Titres de blocs
\newcommand{\exo}[2]{\bigskip\noindent{\large\bfseries Exercice #1}\hfill\textit{#2}\par\nopagebreak\smallskip}
\newcommand{\partie}[1]{\medskip\noindent{\bfseries Partie #1}\par\nopagebreak\smallskip}
\newcommand{\entetesujet}[1]{\begin{center}{\Large\bfseries #1}\end{center}\vspace{0.3em}\hrule\vspace{1.1em}}

% Questions numerotees : niveau 1 = chiffres gras, niveau 2 = a. b. c.
\newlist{questions}{enumerate}{1}
\setlist[questions]{label=\textbf{\arabic*.},leftmargin=1.8em,itemsep=3pt,topsep=3pt}
\newlist{sousq}{enumerate}{1}
\setlist[sousq]{label=\textbf{\alph*.},leftmargin=1.8em,itemsep=2pt,topsep=2pt}


\begin{document}

\entetesujet{Sujet bac 2021 -- Série A}

\exo{1}{8 points}
\partie{A}
On considère les trois fractions suivantes : $a = \dfrac{9}{10}$ ; $b = \dfrac{18}{25}$ et $c = \dfrac{4}{30}$.
\begin{questions}
  \item Simplifier au maximum, s'il y a lieu, chacune des fractions $a$, $b$ et $c$.
  \item Trouver le $\mathrm{PPCM}(10, 25, 15)$.
  \item On rappelle que le plus petit commun dénominateur de deux ou plusieurs fractions est le PPCM des dénominateurs de ces fractions. On considère que le $\mathrm{PPCM}(10, 25, 15)$ est 150.
  \begin{sousq}
    \item Réduire les fractions $\dfrac{9}{10}$, $\dfrac{18}{25}$ et $\dfrac{2}{15}$ au plus petit commun dénominateur.
    \item En déduire le calcul de la somme algébrique $S = \dfrac{9}{10} - \dfrac{18}{25} + \dfrac{4}{30}$. \underline{NB.} On doit détailler les calculs.
  \end{sousq}
\end{questions}

\partie{B}
Soit $P$ la fonction polynôme définie sur $\R$ par $P(x) = 2x^3 + 3x^2 - 11x - 6$.
\begin{questions}
  \item \begin{sousq}
    \item Montrer que $P(x) = (2x + 1)(x^2 + x - 6)$.
    \item Résoudre, dans $\R$, l'équation $x^2 + x - 6 = 0$.
    \item En déduire les solutions de l'équation $P(x) = 0$.
  \end{sousq}
  \item Résoudre, dans $\R$, l'équation $(E) : 2(\ln t)^3 + 3(\ln t)^2 - 11\ln t - 6 = 0$.
\end{questions}

\exo{2}{8 points}
Soit la fonction numérique $g$ de la variable réelle $x$ définie dans $\R^*$ par : $g(x) = \dfrac{x - 1}{x}$. On désigne par $(\mathcal{C}_g)$ la courbe représentative de la fonction $g$ dans un plan muni d'un repère orthonormé $(O,\vi,\vj)$. Unité graphique : 1 cm.
\begin{questions}
  \item \begin{sousq}
    \item Écrire l'ensemble de définition sous forme de réunion d'intervalles.
    \item Écrire $g(x)$ sous forme : $\dfrac{x - 1}{x} = 1 + \dfrac{a}{x}$ où $a$ est un réel que l'on déterminera.
    \item Étudier la limite de $g$ en $-\infty$ et en $0$ à gauche et à droite.
  \end{sousq}
  \item \begin{sousq}
    \item Trouver la fonction dérivée $g'$ de $g$.
    \item Préciser le signe de $g'(x)$ pour tout $x \in \R^*$.
    \item Dresser le tableau de variation de $g$. On donne $\lim\limits_{x\to+\infty} g(x) = 1$.
  \end{sousq}
  \item \begin{sousq}
    \item Trouver les coordonnées des points $A$, $B$ et $C$ d'abscisses respectives $x_A = -1$, $x_B = 1$ et $x_C = 2$.
    \item Représenter dans le plan muni du repère $(O,\vi,\vj)$ la courbe $(\mathcal{C}_g)$ et la droite asymptote horizontale d'équation $y = 1$.
  \end{sousq}
\end{questions}

\exo{3}{4 points}
On considère la série statistique $(X, Y, n_i)$ liée à l'étude du poids $Y_i$ de la larve mesuré en fonction de l'âge $X_i$ de l'insecte. Les résultats de l'étude sont contenus dans le tableau suivant :
\begin{center}\renewcommand{\arraystretch}{1.4}
\begin{tabular}{|c|c|c|c|c|c|c|}
\hline
$X$ (Âge en mois) & 1 & 2 & 3 & 4 & 5 & 6 \\
\hline
$Y$ (Poids en mg) & 1{,}2 & 1{,}8 & 2{,}3 & 2{,}5 & 3 & 3{,}5 \\
\hline
\end{tabular}
\end{center}
\begin{questions}
  \item Calculer les coordonnées du point moyen $G(\overline{X}, \overline{Y})$ du nuage de la série.
  \item \begin{sousq}
    \item Partager le tableau en deux sous-tableaux $T_1$ et $T_2$ tels que dans $T_1$ il y ait les 3 premières valeurs de $X$ et dans $T_2$ les 3 dernières valeurs.
    \item Trouver les points moyens partiels $G_1$ et $G_2$ correspondant respectivement à $T_1$ et $T_2$.
    \item Trouver la droite d'ajustement linéaire de cette série par la méthode de Mayer.
  \end{sousq}
  \item On considère l'équation : $1{,}24x - 3y + 2{,}8 = 0$ représentant une droite d'ajustement de notre nuage. Estimer le poids de la larve au $7^{\text{e}}$ mois si l'évolution se produit dans les mêmes conditions.
\end{questions}

\end{document}
